\section{Fixture Pattern}

The W-series is validated by a small, deliberately shaped fixture family rather
than by claims about any particular election. Four fixture roles recur across
methods and define what a forensic report must and must not say.

\begin{center}
\small
\begin{tabularx}{\textwidth}{lX}
\toprule
Fixture & Role \\
\midrule
No-anomaly & A synthetic package with clean, well-behaved features. A correct method returns low scores and an empty or low-priority queue. \\
Injected outlier & A package with one planted anomalous unit (for example an outlying precinct). The method must surface that unit near the top of the queue. \\
Batch/scanner effect & A package with a systematic batch- or scanner-level deviation. The method must attribute the effect to the right grouping rather than to individual units. \\
False-positive boundary & A package whose pattern is unusual but has an innocent explanation. The report must show why a high score is not proof, exercising the investigation-note and caveat fields. \\
\bottomrule
\end{tabularx}
\end{center}

The fixtures separate two things a forensic method must get right: detecting a
genuinely planted anomaly, and refusing to escalate an unusual-but-explainable
pattern into a fraud claim. The false-positive boundary fixture is the most
important of the four, because it is where the report contract earns its keep:
the same machinery that produces a high score must also produce the caveat that
keeps the score from being read as certification evidence.

These fixtures describe the intended evidence pattern for the track. As with the
report contract, building them is downstream work that depends on normalized
RCOUNT inputs; this paper fixes their roles so that later method papers and any
implementation share one definition of success.
